\documentclass[conference]{IEEEtran}
\IEEEoverridecommandlockouts

\usepackage{cite}
\usepackage{amsmath,amssymb,amsfonts}
\usepackage{algorithm}
\usepackage{algorithmic}
\usepackage{graphicx}
\usepackage{booktabs}
\usepackage{textcomp}
\usepackage{xcolor}
\usepackage[hidelinks]{hyperref}
\usepackage{url}

\def\BibTeX{{\rm B\kern-.05em{\sc i\kern-.025em b}\kern-.08em T\kern-.1667em\lower.7ex\hbox{E}\kern-.125emX}}

% Numbers imported from the analysis pipeline. Regenerated by
% `make analyze`; never edited by hand.
% Auto-generated by analysis/analyze.py -- do not edit by hand.
% Every measurement quoted in the manuscript is defined here.

\newcommand{\AqeCoalMaxTaskMid}{6.2}
\newcommand{\AqeCoalMaxTaskThick}{17.5}
\newcommand{\AqeCoalMaxTaskThin}{1.8}
\newcommand{\AqeCoalSkewMid}{3.5}
\newcommand{\AqeCoalSkewThick}{10.4}
\newcommand{\AqeCoalSkewThin}{1.0}
\newcommand{\AqeCoalTasksMid}{3}
\newcommand{\AqeCoalTasksThick}{3}
\newcommand{\AqeCoalTasksThin}{3}
\newcommand{\AqeCoalWallMid}{6.90}
\newcommand{\AqeCoalWallPreciseMid}{6.90}
\newcommand{\AqeCoalWallPreciseThick}{18.33}
\newcommand{\AqeCoalWallPreciseThin}{2.88}
\newcommand{\AqeCoalWallThick}{18.33}
\newcommand{\AqeCoalWallThin}{2.88}
\newcommand{\AqeNoCoalMaxTaskMid}{5.7}
\newcommand{\AqeNoCoalMaxTaskThick}{16.9}
\newcommand{\AqeNoCoalMaxTaskThin}{1.3}
\newcommand{\AqeNoCoalSkewMid}{30.0}
\newcommand{\AqeNoCoalSkewThick}{105.5}
\newcommand{\AqeNoCoalSkewThin}{6.9}
\newcommand{\AqeNoCoalTasksMid}{16}
\newcommand{\AqeNoCoalTasksThick}{16}
\newcommand{\AqeNoCoalTasksThin}{16}
\newcommand{\AqeNoCoalWallMid}{7.50}
\newcommand{\AqeNoCoalWallPreciseMid}{7.50}
\newcommand{\AqeNoCoalWallPreciseThick}{18.38}
\newcommand{\AqeNoCoalWallPreciseThin}{3.39}
\newcommand{\AqeNoCoalWallThick}{18.38}
\newcommand{\AqeNoCoalWallThin}{3.39}
\newcommand{\AqeSaltShuffleOne}{38}
\newcommand{\AqeSaltShuffleTwo}{38}
\newcommand{\AqeSaltSkewOne}{2.50}
\newcommand{\AqeSaltSkewTwo}{13.91}
\newcommand{\AqeSaltSpeedupOne}{0.87}
\newcommand{\AqeSaltSpeedupTwo}{1.60}
\newcommand{\AqeSaltWallOne}{3.0}
\newcommand{\AqeSaltWallTwo}{11.7}
\newcommand{\AqeShuffleOne}{34}
\newcommand{\AqeShuffleTwo}{34}
\newcommand{\AqeSkewMid}{3.5}
\newcommand{\AqeSkewOne}{1.02}
\newcommand{\AqeSkewThick}{105.5}
\newcommand{\AqeSkewThin}{1.0}
\newcommand{\AqeSkewTwo}{10.39}
\newcommand{\AqeSpeedupMid}{1.01}
\newcommand{\AqeSpeedupOne}{0.91}
\newcommand{\AqeSpeedupThick}{1.02}
\newcommand{\AqeSpeedupThin}{0.91}
\newcommand{\AqeSpeedupTwo}{1.02}
\newcommand{\AqeWallMid}{6.9}
\newcommand{\AqeWallOne}{2.9}
\newcommand{\AqeWallThick}{18.3}
\newcommand{\AqeWallThin}{2.9}
\newcommand{\AqeWallTwo}{18.3}
\newcommand{\BESTK}{2}
\newcommand{\BINDING}{parallelism}
\newcommand{\BaseCoalMaxTaskMid}{5.6}
\newcommand{\BaseCoalMaxTaskThick}{17.0}
\newcommand{\BaseCoalMaxTaskThin}{1.2}
\newcommand{\BaseCoalSkewMid}{36.0}
\newcommand{\BaseCoalSkewThick}{101.0}
\newcommand{\BaseCoalSkewThin}{7.5}
\newcommand{\BaseCoalTasksMid}{16}
\newcommand{\BaseCoalTasksThick}{16}
\newcommand{\BaseCoalTasksThin}{16}
\newcommand{\BaseCoalWallMid}{6.96}
\newcommand{\BaseCoalWallPreciseMid}{6.96}
\newcommand{\BaseCoalWallPreciseThick}{18.62}
\newcommand{\BaseCoalWallPreciseThin}{2.61}
\newcommand{\BaseCoalWallThick}{18.62}
\newcommand{\BaseCoalWallThin}{2.61}
\newcommand{\BaseMaxTaskMid}{5.6}
\newcommand{\BaseMaxTaskThick}{17.0}
\newcommand{\BaseMaxTaskThin}{1.2}
\newcommand{\BaseShuffleOne}{34}
\newcommand{\BaseShuffleTwo}{34}
\newcommand{\BaseSkewMid}{36}
\newcommand{\BaseSkewOne}{7.48}
\newcommand{\BaseSkewThick}{101}
\newcommand{\BaseSkewThin}{7}
\newcommand{\BaseSkewTwo}{101.01}
\newcommand{\BaseSpeedupOne}{1.00}
\newcommand{\BaseSpeedupTwo}{1.00}
\newcommand{\BaseWallMid}{7.0}
\newcommand{\BaseWallOne}{2.6}
\newcommand{\BaseWallThick}{18.6}
\newcommand{\BaseWallThin}{2.6}
\newcommand{\BaseWallTwo}{18.6}
\newcommand{\BcastWallMid}{5.60}
\newcommand{\BcastWallThick}{13.32}
\newcommand{\BcastWallThin}{2.62}
\newcommand{\BtstatMid}{0.06}
\newcommand{\BtstatThick}{1.14}
\newcommand{\BtstatThin}{1.16}
\newcommand{\CORES}{2}
\newcommand{\ControlCrossSessionExcluded}{0}
\newcommand{\ControlMedianPct}{6.5}
\newcommand{\ControlN}{3}
\newcommand{\ControlWorstPct}{18.3}
\newcommand{\DRIVERMEM}{4\,GB}
\newcommand{\GAMMAONE}{0.000561}
\newcommand{\GAMMATWO}{0.0681}
\newcommand{\GammaMid}{inf}
\newcommand{\GammaSpread}{121}
\newcommand{\GammaThick}{0.0681}
\newcommand{\GammaThin}{0.000561}
\newcommand{\HHOT}{445\,743}
\newcommand{\HOTMB}{7.6}
\newcommand{\IdentMid}{no}
\newcommand{\IdentThick}{no}
\newcommand{\IdentThin}{no}
\newcommand{\JAVAVER}{21}
\newcommand{\KEFF}{2}
\newcommand{\KEMPONE}{2}
\newcommand{\KEMPTWO}{2}
\newcommand{\KSTARONE}{nan}
\newcommand{\KSTARTWO}{nan}
\newcommand{\KstarMid}{not identified}
\newcommand{\KstarThick}{not identified}
\newcommand{\KstarThin}{not identified}
\newcommand{\MeasHotMB}{11.41}
\newcommand{\MeasMedMB}{1.35}
\newcommand{\NDIM}{20\,000}
\newcommand{\NFACT}{2\,000\,000}
\newcommand{\NPART}{16}
\newcommand{\NumFits}{3}
\newcommand{\NumGammaEstimable}{2}
\newcommand{\NumIdentified}{0}
\newcommand{\NumNegativeB}{1}
\newcommand{\PHOT}{11}
\newcommand{\PMid}{100}
\newcommand{\PThick}{320}
\newcommand{\PThin}{10}
\newcommand{\PredSlopeRatio}{1818}
\newcommand{\RHOSAT}{3.9}
\newcommand{\RHOVAL}{3.91}
\newcommand{\ROWBYTES}{17}
\newcommand{\RecSlopeRatio}{1818}
\newcommand{\RsqMid}{0.288}
\newcommand{\RsqThick}{0.820}
\newcommand{\RsqThin}{0.312}
\newcommand{\Rsquared}{$R^{2}=0.312$}
\newcommand{\RsquaredTwo}{$R^{2}=0.820$}
\newcommand{\SPARKVER}{3.5.3}
\newcommand{\SaltBestKMid}{2}
\newcommand{\SaltBestKThick}{2}
\newcommand{\SaltBestKThin}{2}
\newcommand{\SaltMaxTaskMid}{3.2}
\newcommand{\SaltMaxTaskThick}{9.1}
\newcommand{\SaltMaxTaskThin}{0.8}
\newcommand{\SaltShuffleOne}{36}
\newcommand{\SaltShuffleTwo}{36}
\newcommand{\SaltSkewOne}{5.82}
\newcommand{\SaltSkewTwo}{68.38}
\newcommand{\SaltSpeedupMid}{1.34}
\newcommand{\SaltSpeedupOne}{0.97}
\newcommand{\SaltSpeedupThick}{1.68}
\newcommand{\SaltSpeedupThin}{0.97}
\newcommand{\SaltSpeedupTwo}{1.68}
\newcommand{\SaltTaskCutMid}{1.7}
\newcommand{\SaltTaskCutThick}{1.9}
\newcommand{\SaltTaskCutThin}{1.6}
\newcommand{\SaltWallMid}{5.2}
\newcommand{\SaltWallOne}{2.7}
\newcommand{\SaltWallThick}{11.1}
\newcommand{\SaltWallThin}{2.7}
\newcommand{\SaltWallTwo}{11.1}
\newcommand{\SelRecR}{1.000}
\newcommand{\SelRecSlope}{11}
\newcommand{\THRESHMB}{8.9}
\newcommand{\UnifRecR}{1.000}
\newcommand{\UnifRecSlope}{20000}
\newcommand{\UnifShuffleOne}{39}
\newcommand{\UnifShuffleTwo}{47}
\newcommand{\UnifSkewOne}{4.97}
\newcommand{\UnifSkewTwo}{17.45}
\newcommand{\UnifSpeedupOne}{0.94}
\newcommand{\UnifSpeedupTwo}{1.59}
\newcommand{\UnifWallOne}{2.8}
\newcommand{\UnifWallTwo}{11.7}


\begin{document}

\title{Adaptive Query Execution or Salting?\\
A Cost Model and Decision Procedure for\\
Skewed Joins in Apache Spark}

\author{\IEEEauthorblockN{Aditya Malviya}
\IEEEauthorblockA{\textit{Independent Researcher} \\
Mumbai, India \\
adityamalviya04692@gmail.com}}

\maketitle

\begin{abstract}
Data skew remains the dominant cause of straggler tasks in distributed joins.
Apache Spark ships an automatic remedy, the skew-join rule of Adaptive Query
Execution (AQE), yet practitioners continue to apply \emph{salting} by hand,
and the two are routinely combined without any account of whether the
combination helps. In the venues we searched we found no
independent, controlled evaluation of Spark's built-in skew-join optimisation,
and no treatment of salting under that name: the technique is transmitted entirely
through vendor documentation and practitioner blogs, none of which prices its
cost. That cost is real. Salting a key across $k$ sub-keys relieves the
straggler in proportion to $1/k$ but replicates the matching rows on the
opposite side in proportion to $k$, so total time is convex in $k$ and has an
interior minimum that the folklore advice to ``pick a reasonable $k$'' cannot
locate.
We contribute three things. First, an analytical cost model
$T(k)=aH/k+bP(k-1)+c$ with the closed-form optimum $k^{*}=\sqrt{aH/(bP)}$,
where $H$ and $P$ are the hot-key row counts on the two sides, together with
two ceilings that bound the useful range: a \emph{saturation} bound
$k\le\rho$, where $\rho$ is the hot key's rows over the mean rows per
partition, beyond which further splitting relieves nothing, and a
\emph{parallelism} bound $k\le C$. Second, a decision procedure that predicts
from data statistics alone whether AQE will suffice, keying on the
one-sided/two-sided distinction that determines whether AQE's
split-and-replicate degenerates. Third, \textsc{skewbench}, an open
benchmark harness that generates Zipfian-skewed workloads with an
aerospace engine-telemetry schema, executes six join strategies, and reports
hardware-independent metrics recovered from the Spark event log.
Measuring 45 configurations on a \NFACT-row join, we find that enabling AQE
produced no improvement distinguishable from measurement noise at any
probe-side thickness (\AqeSpeedupMid$\times$ and \AqeSpeedupThick$\times$
against measurement floors of 0.1\% and 6.5\%), and left the critical path
unchanged: the longest join-stage task went from \BaseMaxTaskThick~s to
\AqeCoalMaxTaskThick~s. Inspecting final adaptive plans shows why --- on a join
whose longest task ran two orders of magnitude beyond the median, Spark's
skew-join rule \emph{never fired}. Neither documented trigger knob was
responsible: sweeping both to their most permissive settings changed nothing,
while lowering \texttt{advisoryPartitionSizeInBytes} below the hot partition
engaged the rule immediately, since that is the size the rule splits into. A
paired control with partition coalescing disabled confirms that AQE's small
residual benefit is coalescing, not skew handling. Selective salting, by
contrast, cut the critical path \SaltTaskCutThick$\times$. Measured in shuffle
records the model's replication term is confirmed exactly --- selective salting
adds \SelRecSlope{} rows per salt value against uniform salting's
\UnifRecSlope, a ratio of \RecSlopeRatio{} matching the predicted $M/P$ at
$R^{2}=1.000$ --- but we decline to claim the closed form for $k^{*}$ is
validated, because the replication coefficient is statistically unidentified in
all \NumFits{} fits. We report that negative result rather than a fitted curve,
and release the cluster configuration under which it becomes testable.
\end{abstract}

\begin{IEEEkeywords}
Apache Spark, data skew, adaptive query execution, salting, cost model,
distributed joins, lakehouse, benchmark
\end{IEEEkeywords}

\section{Introduction}

A join between a large fact table and a smaller dimension table is the most
common heavy operation in an analytics pipeline, and skew in the join key is
the most common reason such a join is slow. When one key carries a large share
of the rows, hash partitioning sends them all to a single reducer, one task
runs far longer than its peers, and the stage finishes when that task finishes.
Adding compute does not help: the straggler is serial. The effect was
characterised for parallel joins three decades ago~\cite{walton1991taxonomy,
dewitt1992skew} and re-confirmed at production scale in every generation of
big-data system since~\cite{chaiken2008scope, ananthanarayanan2010mantri,
kwon2012skewtune}.

Apache Spark~\cite{zaharia2012rdd, armbrust2015sparksql} addresses skew
automatically. Since version 3.0, the \emph{Adaptive Query Execution} (AQE)
framework re-optimises a plan between stages using statistics observed at
runtime, and its \texttt{OptimizeSkewedJoin} rule detects a shuffle partition
that is disproportionately large, splits it, and replicates the matching
partition on the other side so that every split finds its
partners~\cite{databricks_aqe, spark_aqe_docs, xue2024adaptive}.

And yet practitioners keep salting by hand. \emph{Salting} appends a random
suffix drawn from $\{0,\dots,k-1\}$ to the join key on the skewed side and
replicates the matching rows $k$ times on the other side, so that one heavy key
becomes $k$ lighter ones. The technique is ubiquitous in industry and, as we
establish in Section~\ref{sec:background}, entirely absent from the
peer-reviewed literature under that name. It is transmitted through vendor
documentation, conference talks and blog posts, none of which prices what it
costs.

That silence matters, because salting is not free. Every additional salt value
multiplies the rows that must be replicated and shuffled on the opposite side.
The straggler term falls as $1/k$; the replication term rises linearly in $k$.
Total time is therefore \emph{convex} in $k$ with an interior minimum, and the
standard advice --- ``pick a reasonable $k$, say 16 or 32'' --- is not advice at
all. It names a point on a curve whose shape nobody has published.

\subsection{The questions this paper answers}

\begin{enumerate}
\item Given a join's key distribution and Spark's configuration, will AQE's
skew-join rule suffice on its own, or is explicit salting still required?
\item If salting is required, what value of $k$ minimises total cost, and can
it be derived from statistics available before the query runs rather than found
by trial and error?
\item Does combining AQE with salting help, hurt, or do nothing?
\end{enumerate}

\subsection{Contributions}

\begin{itemize}
\item \textbf{A cost model for salt cardinality.} We formulate the salting
trade-off as $T(k) = aH/k + bP(k-1) + c$, where $H$ and $P$ are the hot-key row
counts on the skewed and probe sides, and derive the closed-form optimum
$k^{*}=\sqrt{aH/(bP)}$. The single free ratio $\gamma = a/b$ is a property of
the engine and hardware, calibrated once and reused
(Section~\ref{sec:costmodel}).

\item \textbf{Two ceilings that bound the useful range.} Beyond
$k=\rho$, where $\rho$ is the hot key's rows divided by the mean rows per
partition, the hot task is no longer the straggler and further splitting
relieves nothing; beyond $k=C$, the number of concurrently schedulable task
slots, the sub-partitions queue rather than overlap. \emph{Salting beyond
$\min(\rho, C)$ is always waste} --- a bound that follows directly from the
model and that we could not find stated anywhere in the practitioner
literature.

\item \textbf{A decision procedure.} A rule that predicts, from data statistics
and Spark configuration alone, whether AQE will trigger, whether it will
suffice, and what $k$ to use otherwise. Its hinge is the distinction between
one-sided skew, which AQE is designed for, and two-sided skew, in which
split-and-replicate degenerates because each split inherits a large partner
(Section~\ref{sec:decision}).

\item \textbf{\textsc{skewbench}, an open harness.} A reproducible benchmark
that generates Zipfian-skewed workloads with an aerospace engine-telemetry
schema, executes six join strategies, and recovers hardware-independent metrics
--- shuffle bytes, spill bytes, and the within-stage ratio of maximum to median
task duration --- by parsing the Spark event log. It runs on a commodity laptop
(Section~\ref{sec:methodology}).

\item \textbf{An empirical evaluation} over 48 configurations, which confirms
the decay of AQE's benefit as the probe side thickens, identifies the
precondition that actually blocked Spark's skew-join rule from ever firing, and
quantifies the replication-cost gap between the two salting variants. It also
reports, as a negative result, that the convexity of $T(k)$ does not resolve at
single-node scale for selective salting (Section~\ref{sec:results}).
\end{itemize}

\subsection{Scope, stated plainly}

This is an evaluation and decision-procedure paper, not an algorithm paper. We
propose no new partitioner; the space of skew-resilient partitioning algorithms
is already well populated~\cite{irandoost2019skewslr, gufler2011skew,
myung2016handling, shen2020srspark, tang2021intermediate}. Our claim is that
the prior question --- \emph{which of the mechanisms that already ship should a
practitioner use, and how should it be parameterised} --- has gone unanswered,
and that it is answerable.

All measurements were taken on a single node. We are explicit about what that
does and does not license. Absolute wall-clock times do not transfer to a
cluster and we make no claim that they do; shuffle volume, spill volume and
within-stage task-time ratios are properties of the physical plan and the data
partitioning rather than of the machine, and it is on those that the
quantitative claims rest. Section~\ref{sec:threats} treats this at length.

\section{Background and Related Work}
\label{sec:background}

\subsection{Skew in parallel joins}

Skew in parallel joins is old and well understood in the abstract. Walton et
al.~\cite{walton1991taxonomy} gave a taxonomy and performance model of skew
effects; DeWitt et al.~\cite{dewitt1992skew} proposed practical handling
including range partitioning with subset replication for heavy keys --- the
direct ancestor of what practitioners now call salting. In the MapReduce era,
SkewTune~\cite{kwon2012skewtune, kwon2012skewtunedemo} repartitioned a
straggler's remaining input at runtime, and Mantri~\cite{ananthanarayanan2010mantri}
established at production scale that outlier tasks dominate job latency.

A substantial line of work proposes new partitioners. For Spark specifically:
intermediate-data partitioning for skew mitigation~\cite{tang2021intermediate}
and skew-resilient adaptive parallel processing~\cite{shen2020srspark}. In the
MapReduce lineage from which these descend: sampling-based and cost-aware
partitioning schemes~\cite{myung2016handling, gufler2011skew}. Irandoost et
al.~\cite{irandoost2019skewslr} survey the MapReduce field systematically.
Their review discusses neither Adaptive Query Execution nor salting --- fairly
enough for AQE, which shipped in Spark 3.0 after that review was written, but
the omission is indicative: the literature has been optimising the mechanism
while the question of \emph{selecting} among deployed mechanisms went unasked.

Metwally's work on scaling and load-balancing equi-joins~\cite{metwally2025amjoin}
is the closest formal treatment of the idea salting embodies. AM-Join combines
randomised key expansion with replication on the opposite side and analyses
doubly-skewed keys as the hard case. It is the right formal ancestor to cite,
and it is not a study of the technique practitioners actually run.

\subsection{Adaptive Query Execution}

AQE re-optimises between stages using runtime statistics: coalescing
post-shuffle partitions, converting sort-merge joins to broadcast joins when
the observed side turns out small, and splitting skewed
partitions~\cite{databricks_aqe, spark_aqe_docs}. Xue et
al.~\cite{xue2024adaptive} describe the production descendant of this machinery
at lakehouse scale and state its skew handling explicitly: the rule discovers
skew on the join keys of a shuffled hash join, ``which manifests as a few
partitions containing significantly more data than others,'' and eliminates it
by ``logically splitting those large, consumer-side partitions into smaller,
more balanced partitions.''

That paper is peer-reviewed, and it is the authoritative account of the
mechanism. It is also a systems-description paper written by the vendor and
evaluated on the vendor's engine. What it does not provide --- and what does not
exist elsewhere --- is an independent, controlled evaluation of the open-source
rule against the alternative practitioners actually reach for.

Adjacent work confirms the shape of the gap rather than filling it. Learned and
adaptive Spark optimisers --- AQORA~\cite{he2025aqora},
LOCAT~\cite{xin2022locat}, fine-grained parameter tuning~\cite{lyu2024sparkopt},
zero-execution configuration tuning~\cite{suri2025zeroexec} --- address join
ordering, broadcast selection and configuration knobs. AQORA operates
\emph{inside} the AQE framework and does not address skew. Zhang et
al.~\cite{zhang2023adaptive} evaluate adaptive query processing rigorously, but
on single-node relational engines rather than Spark.

\subsection{The specific absence}

We searched dblp, the arXiv full-metadata API, Crossref, and domain-scoped
searches over the ACM Digital Library, IEEE Xplore, the VLDB proceedings,
Springer, ScienceDirect and MDPI. Two negative results are load-bearing for
this paper, and we state them with their limits.

\textbf{No independent evaluation of Spark's skew-join rule.} We found no
peer-reviewed paper whose contribution is an evaluation of
\texttt{spark.sql.adaptive.skewJoin}. The nearest comparative study, Phan et
al.~\cite{phan2022skewjoin}, compares hash-based, range-based and
multi-dimensional range partitioning on MapReduce and Spark --- and does not
consider AQE. The mechanism is otherwise documented only in vendor
material~\cite{databricks_aqe, spark_aqe_docs} and described from the
provider's perspective~\cite{xue2024adaptive}.

\textbf{No peer-reviewed treatment of salting.} A full-metadata arXiv search
for \texttt{abs:salting AND abs:join} returns five papers, all in
condensed-matter physics and astronomy. Searches for ``key salting'' and
``salted join'' across the major publisher databases return only unrelated
hash-join work. Every source describing the technique as practitioners apply it
is a blog post, a tutorial, or vendor documentation. Its formal analogue is
randomised key expansion with replication~\cite{metwally2025amjoin,
dewitt1992skew}, never under that name and never as the object of study.

We note the limits of this claim. Semantic Scholar and OpenAlex were
unavailable during our search, so recall rests on dblp, arXiv, Crossref and
domain-scoped web search, and dblp's coverage of some journals is imperfect ---
Phan et al. is itself not indexed there. We therefore claim these absences as
strong evidence rather than as proof.

\subsection{Cost models and benchmarking practice}

Cost modelling for big-data query processing is well
developed~\cite{siddiqui2020costmodels, zhou2010scopeopt, leis2021costoptimal},
including for shuffle specifically~\cite{shen2020magnet}, but we are not aware
of a cost model for the salting decision. Methodologically we follow
LST-Bench~\cite{camacho2024lstbench} in building the benchmark from composable,
declared workload patterns, and Nurdin et al.~\cite{nurdin2026lakehouseperf} in
treating runtime variance as a first-class threat: they measure roughly
twofold run-to-run variation for identical lakehouse queries on shared cloud
compute, which is precisely why our timed measurements are taken on dedicated
local compute and our transferable claims rest on deterministic counters.

\section{Problem Formulation}
\label{sec:problem}

\subsection{Notation}

Consider an inner equi-join $F \bowtie_{\kappa} D$ on key $\kappa$, where $F$ is
the larger, skewed side with $N$ rows and $D$ the probe side with $M$ rows.
Spark hash-partitions both sides into $p$ shuffle partitions
(\texttt{spark.sql.shuffle.partitions}) and executes a sort-merge join with $C$
concurrently schedulable task slots.

Let the key frequencies on $F$ follow a truncated Zipf law over $n$ distinct
keys, $\Pr[\kappa = i] \propto i^{-\theta}$, and let key $1$ be the hot key.
Write

\begin{itemize}
\item $H$ for the rows of $F$ carrying the hot key,
\item $P$ for the rows of $D$ carrying the hot key.
\end{itemize}

Both are obtainable before execution, from table statistics, a sampled
frequency sketch, or the catalogue.

It is convenient to express $P$ relative to the probe side's mean rows per key,
$\bar{m} = M/n$, through a \emph{hot-key multiplier}
\begin{equation}
\mu \;=\; \frac{P}{\bar{m}} ,
\label{eq:mu}
\end{equation}
so that $\mu = 1$ is a probe side with no concentration at all --- one-sided
skew --- and larger $\mu$ makes the skew progressively two-sided. Treating the
degree of two-sidedness as continuous rather than binary matters for two
reasons, one practical and one methodological, and we return to both in
Section~\ref{sec:methodology}.

\subsection{The skew ratio}

The quantity that governs everything is dimensionless:

\begin{equation}
\rho \;=\; \frac{H}{N/p}
\label{eq:rho}
\end{equation}

the hot key's rows divided by the mean rows per shuffle partition. It says how
many ordinary partitions' worth of work has been concentrated into one.

Equation~\eqref{eq:rho} is a property of the data and the partition count
alone. It contains no hardware term, which is what allows a rule built on it to
be stated once and applied on any cluster. When $\rho \le 1$ there is no
straggler to remove. When $\rho \gg 1$ the stage's completion time is
essentially the hot task's completion time, regardless of how many executors are
provisioned.

\subsection{Why $\mu$ cannot be taken to its limit}

There is a natural temptation to construct the two-sided case by drawing both
sides from the same Zipf law. It does not work, and the reason is worth stating
because it constrains any experiment in this area.

If both sides follow the same law, the hot key contributes $H \times P$ rows to
the join output on its own. At the scales used here that is order $10^{9}$
rows from a single key, a thousandfold amplification of the fact table --- the join has become a near-cartesian product, and
any measurement of it reports the cost of materialising that product rather
than the cost of the skew mitigation under test. Bounding $\mu$ explicitly
keeps the output tractable and, more usefully, promotes $P$ to an independent
variable. Equation~\eqref{eq:kstar} predicts that the optimal salt cardinality
falls as $1/\sqrt{P}$; that prediction is only testable if $P$ can be swept.

\subsection{Two mechanisms, one problem}

\subsubsection{AQE skew-join}
Spark's \texttt{OptimizeSkewedJoin} rule classifies a shuffle partition as
skewed when its size exceeds \emph{both} an absolute threshold
(\texttt{skewedPartitionThresholdInBytes}) and a multiple
(\texttt{skewedPartitionFactor}) of the median partition size. It then splits
the offending partition into sub-partitions of roughly the advisory size and
\emph{replicates the corresponding partition of the other side} to each split,
so that every sub-partition retains its join partners~\cite{spark_aqe_docs,
xue2024adaptive}.

The replication is the hinge. If the other side's matching partition is small
--- the one-sided case --- replication is nearly free and the rule works as
advertised. If the other side is also concentrated on the same key --- the
two-sided case --- then each of the $s$ splits inherits a large partner, the
probe-side work multiplies by $s$, and the relief on the critical path is
partly or wholly cancelled. Metwally~\cite{metwally2025amjoin} identifies
doubly-skewed keys as the hard case for exactly this reason.

\subsubsection{Salting}
Salting rewrites the key. On $F$, the hot key is replaced by
$(\kappa, \sigma)$ with $\sigma$ drawn uniformly from $\{0,\dots,k-1\}$; on
$D$, each matching row is emitted $k$ times, once per salt value. One heavy
key becomes $k$ lighter ones.

We distinguish two variants, because their costs differ by orders of magnitude
and the literature does not separate them:

\begin{description}
\item[Uniform salting] every row of $D$ is replicated $k$ times. This is the
form shown in most tutorials. Replication cost scales with $M$.
\item[Selective salting] only rows whose key is hot are replicated; all other
rows carry salt $0$ on both sides and join exactly as before. Replication cost
scales with $P$.
\end{description}

Since $P \ll M$ in any realistic fact--dimension join, the distinction is not a
detail. We evaluate both.

\subsection{What is actually being traded}

Salting does not remove work; it moves it. Splitting the hot key across $k$
sub-keys reduces the critical task's input from $H$ to approximately $H/k$
rows, and simultaneously introduces $P(k-1)$ probe-side rows that did not
previously exist. One term falls, the other rises, and the practitioner is
choosing a point on that curve without being told there is a curve.

The next section prices it.

\section{A Cost Model for Salt Cardinality}
\label{sec:costmodel}

\subsection{Formulation}

Under selective salting with cardinality $k$, total stage time decomposes into
three parts.

\emph{The critical task.} The hot key's $H$ rows are spread over $k$
sub-partitions, so the heaviest task processes approximately $H/k$ rows. Write
$a$ for the marginal cost of a row on the critical path.

\emph{Replication.} The $P$ matching probe-side rows are emitted $k$ times
instead of once, adding $P(k-1)$ rows to shuffle, sort and merge. Write $b$ for
the marginal cost of a replicated row.

\emph{Everything else.} Scanning both inputs, joining the cold keys, and the
terminal aggregation are all independent of $k$. Collect them into $c$.

\begin{equation}
T(k) \;=\; a\,\frac{H}{k} \;+\; b\,P\,(k-1) \;+\; c ,
\qquad k \ge 1 .
\label{eq:model}
\end{equation}

\subsection{The optimum}

$T$ is the sum of a convex decreasing term and a linear increasing term, hence
strictly convex on $k>0$, with a unique interior minimum. Differentiating,

\begin{equation}
\frac{dT}{dk} \;=\; -\frac{aH}{k^{2}} \;+\; bP \;=\; 0
\qquad\Longrightarrow\qquad
\boxed{\;k^{*} \;=\; \sqrt{\frac{aH}{bP}} \;=\; \sqrt{\gamma\,\frac{H}{P}}\;}
\label{eq:kstar}
\end{equation}

where $\gamma = a/b$ is the ratio of the two marginal costs. Four properties of
Equation~\eqref{eq:kstar} are worth drawing out.

First, \textbf{it depends on the workload only through $H/P$}, the ratio of
hot-key rows on the two sides. That ratio is cheap to estimate and is exactly
what a frequency sketch on the join key already yields.

Second, \textbf{$\gamma$ ought to be a property of the engine and the hardware
rather than of the query}: row-processing cost on the critical path and
replicated-row shuffle cost both scale with the machine, so their ratio should
be more stable than either alone, and one calibration should serve many
workloads. We state this as a hypothesis rather than a result because our own
measurements do not support it --- see Section~\ref{sec:results}, where fitted
$\gamma$ varies by \GammaSpread$\times$ across three workloads on one machine.
That instability is diagnostic of the fits being dominated by the $k=1$ to
$k=2$ step rather than by the replication term, and it is the second reason,
alongside the unresolved convexity, that we do not claim the closed form is
validated.

Third, \textbf{$k^{*}$ falls as $1/\sqrt{P}$}. Thickening the probe side's hot
key makes every additional salt value more expensive while leaving straggler
relief unchanged, so the optimum retreats. This is the model's sharpest
falsifiable prediction, and Section~\ref{sec:results} tests it directly by
sweeping $\mu$ over a factor of 32.

Fourth, and least intuitively, \textbf{$k^{*}$ grows only as the square root of
the skew}. Doubling the hot key's dominance does not double the salt
cardinality; it multiplies it by $\sqrt{2}$. This is the structural reason
aggressive salting is so often counterproductive: intuition says respond to
severe skew with large $k$, and the mathematics says respond with a modestly
larger one.

\subsection{Two ceilings}

Equation~\eqref{eq:kstar} is unconstrained, and taken literally it will
sometimes recommend a $k$ that buys nothing at all.

\subsubsection{Saturation}
Once the hot partition has been divided below the size of an ordinary
partition, the hot task stops being the straggler --- some other partition now
finishes last --- and further division relieves nothing while continuing to pay
replication. From Equation~\eqref{eq:rho}, that point is
\begin{equation}
k_{\text{sat}} \;=\; \rho \;=\; \frac{H}{N/p} .
\end{equation}
\textbf{Salting beyond $\rho$ is always waste.} The bound is immediate from the
model, requires no calibration, and we have not found it stated in the
practitioner literature.

\subsubsection{Parallelism}
Sub-partitions beyond the number of concurrently schedulable task slots $C$
queue rather than overlap, so the straggler term stops falling at $k=C$ while
the replication term keeps rising:
\begin{equation}
k_{\text{par}} \;=\; C .
\end{equation}

\subsubsection{The operational rule}
\begin{equation}
k_{\text{eff}} \;=\; \max\Bigl(1,\;\bigl\lfloor \min(k^{*},\,\rho,\,C) \bigr\rceil\Bigr) .
\label{eq:keff}
\end{equation}

Which of the three binds is itself diagnostic. If $k^{*}$ binds, the workload
is replication-limited and the probe side is the thing to shrink. If $\rho$
binds, the skew is mild relative to the partition count and raising $p$ would
serve better than salting. If $C$ binds, the cluster is the limit and salting
further is pure overhead --- the common case on small clusters, and the one in
which tutorial defaults do the most damage.

\subsection{Uniform salting}

For uniform salting the replication term scales with the whole probe table, so
Equation~\eqref{eq:model} becomes $T_{u}(k) = aH/k + bM(k-1) + c$ and
$k_{u}^{*} = \sqrt{\gamma H/M}$. Since $M \gg P$, the optimum is
correspondingly smaller --- often below $2$, meaning uniform salting is
frequently not worth doing at all. This falls out of the model rather than
having to be discovered empirically, and it is a good illustration of why
pricing the technique is worth the trouble.

\subsection{Calibration}

Given measurements $\{(k_i, T_i)\}$, Equation~\eqref{eq:model} is \emph{linear}
in $(a,b,c)$ under the basis $[H/k,\; P(k-1),\; 1]$. Fitting is therefore
ordinary linear least squares: no initial guess, no local minima, no optimiser
hyperparameters, and a fit that is a deterministic function of the
measurements. Three distinct values of $k$ suffice to identify three
coefficients; we use seven.

\section{The Decision Procedure}
\label{sec:decision}

Section~\ref{sec:costmodel} answers ``how much salt.'' The prior question is
whether to salt at all, given that AQE may already handle the case. The
procedure in Algorithm~\ref{alg:decide} answers both, from statistics available
before execution.

\begin{algorithm}[t]
\caption{Skew mitigation decision procedure. Branch structure and thresholds
are a direct transcription of \texttt{costmodel.decide} and
\texttt{costmodel.recommend\_k} in the released artifact.}
\label{alg:decide}
\begin{algorithmic}[1]
\REQUIRE $H$, $P$, $N$, $M$, $n$, row width $w$, partitions $p$, slots $C$;
         AQE parameters $\beta$ (threshold), $\phi$ (factor), $\alpha$
         (advisory size); optional calibrated $\gamma$
\ENSURE  a strategy, and a salt cardinality where one is required
\STATE $\rho \leftarrow H p / N$
\IF{$\rho < \rho_{\min}$}
    \RETURN \textsc{DoNothing} \COMMENT{no straggler exists}
\ENDIF
\STATE $B_{\text{hot}} \leftarrow H w$;\quad
       $B_{\text{med}} \leftarrow (N-H)\,w/(p-1)$;\quad
       $\Theta \leftarrow \max(\beta,\ \phi B_{\text{med}})$
\IF{$B_{\text{hot}} \le \alpha$}
    \RETURN \textsc{Salt} \COMMENT{no split possible; the rule is a no-op}
\ENDIF
\IF{$B_{\text{hot}} \le \Theta$}
    \RETURN \textsc{Salt} \COMMENT{AQE will not classify the partition as skewed}
\ENDIF
\STATE $\mu \leftarrow P n / M$
\IF{$\mu \le \mu_{\min}$}
    \RETURN \textsc{AqeOnly}
\ELSE
    \RETURN \textsc{Aqe}$+$\textsc{Salt}
\ENDIF
\medskip
\STATE \textbf{if a salt cardinality is required:}
\IF{$\gamma$ is calibrated \textbf{and} the fit is identified}
    \STATE $k^{*} \leftarrow \sqrt{\gamma H / P}$
\ELSE
    \STATE $k^{*} \leftarrow \textsc{Undefined}$ \COMMENT{omit from the minimum}
\ENDIF
\RETURN $k_{\text{eff}} \leftarrow \max\bigl(1, \lfloor\min(k^{*}, \rho, C)\rceil\bigr)$
\end{algorithmic}
\end{algorithm}

The procedure has three branches, and each corresponds to a distinct failure or
success mode.

\subsection{Branch 1: no straggler}
When $\rho < \rho_{\min}$ (we use $\rho_{\min}=2$) the hot key amounts to less
than a couple of ordinary partitions. There is nothing to relieve, and both
mechanisms would add overhead for no gain. This branch exists because
practitioners routinely salt joins that were never skewed enough to matter,
having diagnosed ``skew'' from a slow job rather than from the key
distribution.

\subsection{Branch 2: AQE will not fire}
AQE's rule is conjunctive in \emph{three} conditions, not the two that are
usually discussed. The partition must exceed an absolute byte threshold, a
multiple of the median, \emph{and} the advisory partition size --- because that
last is the size the rule splits into, so a hot partition smaller than it
cannot be divided at all and the optimisation is a no-op however skewed it is.
Section~\ref{sec:results} finds this third condition to be the binding one in
practice, and it is the one practitioners do not tune. A join can be badly skewed in the sense
that matters --- one task running many times longer than its peers --- while
sitting below the absolute threshold, whose default is large. In that regime
AQE does nothing at all, silently, and the practitioner who has enabled it
believes the problem is handled. The remedy is either to salt explicitly or to
lower the threshold; the procedure reports which, rather than leaving the user
to discover the non-event from the Spark UI.

\subsection{Branch 3: the one-sided/two-sided hinge}
This is the substantive branch. When skew is one-sided, splitting the fact-side
partition and replicating the small matching probe-side partition to each split
is close to free, and AQE resolves the straggler. When both sides carry the
same hot key, each of the $s$ splits inherits a large partner. Probe-side work
grows with $s$, and what looked like $s$-fold relief on the critical path is
substantially cancelled.

Detecting which case obtains is cheap, and it is exactly
Equation~\eqref{eq:mu}: compare $P$ against the probe side's mean rows per key.
We use $\mu_{\min}=2$ throughout, in the
text and in the implementation. This is a single group-by on a sample, and it is the
highest-value statistic a practitioner can collect before deciding how to fix a
slow join --- it selects the mechanism \emph{and}, through
Equation~\eqref{eq:kstar}, sets its parameter.

\subsection{Worked example}

Take the configuration measured in Section~\ref{sec:results}: $N=\NFACT$ rows,
$p=\NPART$, $H\approx\HHOT$ rows, $P\approx\PHOT$ rows.

The procedure needs the size of the hot shuffle partition, and here the
analytical estimate and the measurement disagree in a way worth stating.
Estimating $B_{\text{hot}} = Hw$ gives \HOTMB~MB; the event log records the
largest join-stage shuffle partition at \MeasHotMB~MB. The estimate is low
because the hot partition also absorbs whatever cold keys hash to it. In the
other direction, $B_{\text{med}}$ estimates the median at \THRESHMB~MB against
a measured \MeasMedMB~MB, because it ignores variance among the cold
partitions. $Hw$ and $(N-H)w/(p-1)$ are therefore best read as a lower and an
upper bound respectively, and where an event log is available the measured
sizes should be used instead --- as they are in Section~\ref{sec:results}.

Evaluated on the measured sizes, with $\rho=\RHOVAL$ a straggler exists; the
hot partition clears both documented trigger conditions; and it is nonetheless
smaller than the advisory split size, so the procedure returns \textsc{Salt}
with the reason that no split is possible. Section~\ref{sec:results} reports
what Spark actually did.

\section{Experimental Methodology}
\label{sec:methodology}

\subsection{\textsc{skewbench}}

We implement the study as an open harness rather than a set of scripts, because
the artifact is as much the contribution as the numbers. \textsc{skewbench}
comprises a workload generator, six join strategies, an event-log metrics
parser, an experiment driver, and the cost model of
Section~\ref{sec:costmodel}. Every result in this paper is produced by a single
command and every figure and table is generated from the results file with no
manual transcription.

\subsection{Workload generation}

The generator produces two Parquet tables with an aerospace engine-telemetry
schema, following the NASA C-MAPSS turbofan degradation
convention~\cite{saxena2008cmapss} --- unit identifier, operational cycle,
three operational settings, and $N_s$ sensor channels --- so that the synthetic
generator can be exchanged for the real dataset without touching the benchmark.

\texttt{fact\_telemetry} is the skewed side: one row per sensor snapshot, with
\texttt{engine\_id} drawn from a truncated Zipf law over $n$ engines. This
mirrors a real fleet, in which a small number of problem units generate a
disproportionate share of recorded events. \texttt{dim\_maintenance} is the
probe side: one row per maintenance work order. Its keys are uniform except for
the hot key, which carries $\mu$ times the mean per-key row count.

The choice to parameterise two-sidedness by $\mu$ rather than by drawing the
probe side from the same Zipf law is the single most consequential design
decision in the harness, and it is not cosmetic. Under a doubly-Zipf design the
hot key alone yields $H \times P$ output rows. With $H \approx 4.4\times10^{5}$
and $P \approx 4.4\times10^{3}$ that is roughly $2\times10^{9}$ rows from a
single key --- a thousandfold amplification of a two-million-row fact table,
which we observed as a run that made no progress and abandoned. Such a workload
does not measure skew mitigation; it measures the cost of materialising a
near-cartesian product, and every strategy looks equally bad because all of
them are dominated by output volume. Bounding $\mu$ keeps the
output tractable and promotes $P$ from an artefact of the generator to a swept
independent variable, which is what makes the $1/\sqrt{P}$ prediction testable
at all.

Skewed variants of TPC-H exist~\cite{microsoft_skewgen, tpch} and the harness
accepts them, but they do not provide a controllable $\mu$, which is the
condition our decision rule turns on. Generation is deterministic given a seed,
and the manifest written alongside the data records $H$ and $P$ exactly, so no
later stage must re-derive them.

\subsection{Strategies evaluated}

\begin{enumerate}
\item \textbf{baseline} --- sort-merge join, AQE disabled, no mitigation.
\item \textbf{aqe} --- AQE enabled with skew-join optimisation.
\item \textbf{salt\_uniform}$(k)$ --- whole probe side replicated $k$ times,
AQE disabled.
\item \textbf{salt\_selective}$(k)$ --- hot-key rows only replicated, AQE
disabled.
\item \textbf{aqe\_salt}$(k)$ --- AQE enabled \emph{and} selective salting.
\item \textbf{broadcast} --- broadcast hash join, where the probe side fits.
\end{enumerate}

Automatic broadcast is disabled throughout
(\texttt{autoBroadcastJoinThreshold} $=-1$) except in the broadcast arm, so
that a strategy under test is never silently replaced by a different plan. Each
arm's executed physical plan is captured into the results file, so that the
claim ``this was a sort-merge join'' is auditable rather than asserted.

\subsection{Metrics}

We separate metrics by what they license.

\textbf{Deterministic.} Shuffle bytes read and written, records shuffled, and
bytes spilled to memory and disk are properties of the physical plan and the
partitioning. They are identical on a laptop and on a hundred-node cluster
given the same plan and the same data, and the quantitative claims about
replication cost rest on them.

\textbf{Structural.} The \emph{skew ratio} --- maximum over median task
duration within the join stage --- and the task-time coefficient of variation
are within-run ratios. They cancel most machine-specific scaling while still
capturing the straggler that skew produces, and they are how we measure whether
a mitigation actually worked as opposed to merely running faster.

\textbf{Machine-dependent.} Wall-clock and executor run time. Reported, used to
calibrate $\gamma$, and never load-bearing for a transferable claim.

Metrics are recovered by parsing the Spark event log --- a JSON-lines file
written by Spark itself --- rather than through a listener callback. The log
survives the driver exiting and can be re-analysed later without re-running
anything, which is what makes a result auditable months after the fact. Each
measurement is tagged with a unique Spark job group, giving an exact rather
than timestamp-inferred mapping from a results row to the tasks that produced
it.

\subsection{Execution protocol}

Each cell runs one warmup execution, discarded, followed by five timed
repetitions. Warmup matters because the first execution of a plan pays JIT
compilation, class loading and page-cache misses, and including it inflates the
mean unevenly across arms. We report the \emph{median} and interquartile range
rather than mean and standard deviation, because run times on a shared machine
have a long right tail.

Materialisation is forced with Spark's \texttt{noop} sink, so no rows are
written to disk or shipped to the driver and the measurement is of the join and
its aggregation alone. The terminal operation is a grouped aggregation touching
columns from both sides, not a \texttt{count()}: counting invites the optimizer
to prune the payload and, on some plans, to skip the probe entirely, which
would measure column pruning rather than the strategy under test.

\subsection{Platform}

Measurements were taken on a single node with \CORES\ cores and \DRIVERMEM\ of
driver memory, Apache Spark \SPARKVER\ on OpenJDK \JAVAVER, with dedicated
local compute --- not shared serverless. This choice is deliberate: Nurdin et
al.~\cite{nurdin2026lakehouseperf} measure roughly twofold run-to-run variance
for identical lakehouse queries on shared cloud compute, which would make
timing claims meaningless. We are trading absolute scale for measurement
control, and Section~\ref{sec:threats} states what that costs us.

The full configuration grid, the raw event logs, and the exact software
versions are included in the artifact.

\section{Results}
\label{sec:results}

We report 45 configurations: three probe-side thicknesses
($P \in \{\PThin, \PMid, \PThick\}$ rows on the hot key) crossed with six
strategies, a salt sweep $k \in \{1,\dots,32\}$, and --- for the AQE arms --- a
paired run with partition coalescing disabled. Each cell runs one discarded
warmup and five timed repetitions, in randomised order. Every number below is
generated from the results file by the released analysis code; none is
transcribed by hand.

\subsection{The measurement floor}

Selective salting at $k=1$ compiles to exactly the same physical plan as the
unmitigated baseline, so the difference between them is pure measurement noise.
We report it first, because it bounds what any other number in this section can
mean:

\begin{center}
\begin{tabular}{lrrr}
\toprule
$P$ & \PThin & \PMid & \PThick \\
\midrule
Baseline wall (s)                & \BaseWallThin & \BaseWallMid & \BaseWallThick \\
Identical-plan control (s)       & 3.09 & 6.97 & 19.84 \\
\textbf{Measurement floor}       & \textbf{18.3\%} & \textbf{0.1\%} & \textbf{6.5\%} \\
\bottomrule
\end{tabular}
\end{center}

The floor varies sharply with run length, and this governs how the rest of the
section may be read. At $P=\PMid$ the two identical plans agree to 0.1\%, and a
difference of a few percent is meaningful. At $P=\PThin$ they differ by 18.3\%:
those runs last under three seconds, short enough that JVM and scheduler noise
dominates, and \textbf{no claim can be supported at that probe thickness at
all}. We report the $P=\PThin$ column throughout for completeness and draw no
conclusions from it.

Two points about how this number is produced. All 45 cells ran in a single
Spark application: an earlier attempt gap-filled four cells from a second
application, and those cells were systematically slower and drifting within
their own run, which inflated this control from 2\% to 29\%. Wall-clock
comparisons across Spark applications measure JVM warmth, not plans, and the
harness now records the application id and refuses to compare across it.
Separately, an earlier version executed cells in configuration order, placing
the baseline first on a cold JVM; that inflated the same control to 17.7\% and,
with it, every speedup measured against that baseline. Cell order is now
randomised under a fixed seed.

\subsection{The baseline is severely skewed}

Without mitigation, the join stage's longest task takes \BaseMaxTaskThin~s,
\BaseMaxTaskMid~s and \BaseMaxTaskThick~s at the three probe thicknesses, against
median tasks of a fraction of a second. At $P=\PThick$ a single task runs
roughly two orders of magnitude longer than its median peer. This is a textbook
straggler, and no amount of additional compute would remove it.

\subsection{AQE did not help at any probe thickness}

This is the paper's central measurement, and it is a null result.

\begin{center}
\begin{tabular}{lrrr}
\toprule
$P$ & \PThin & \PMid & \PThick \\
\midrule
Baseline wall (s)            & \BaseWallThin & \BaseWallMid & \BaseWallThick \\
AQE wall (s)                 & \AqeWallThin & \AqeWallMid & \AqeWallThick \\
\textbf{AQE speedup}         & \textbf{\AqeSpeedupThin$\times$} & \textbf{\AqeSpeedupMid$\times$} & \textbf{\AqeSpeedupThick$\times$} \\
\midrule
Baseline longest task (s)    & \BaseMaxTaskThin & \BaseMaxTaskMid & \BaseMaxTaskThick \\
AQE longest task (s)         & \AqeCoalMaxTaskThin & \AqeCoalMaxTaskMid & \AqeCoalMaxTaskThick \\
\bottomrule
\end{tabular}
\end{center}

At $P=\PMid$, where the measurement floor is 0.1\%, AQE returns
\AqeSpeedupMid$\times$. At $P=\PThick$, floor 6.5\%, it returns
\AqeSpeedupThick$\times$ --- inside the floor. And the row that settles it is
the last one: enabling Adaptive Query Execution on a join whose longest task
ran \BaseMaxTaskThick~s produced a longest task of \AqeCoalMaxTaskThick~s.
\textbf{The critical path is not shortened at any probe thickness; in all three
it is marginally longer.}

We are careful about what this does and does not say. It is not a claim that
AQE is useless in general --- it does other useful things, and one of them
shows up below. It is a claim that on \emph{this} join, with Spark's skew
machinery enabled and a straggler two orders of magnitude out, AQE delivered no
improvement distinguishable from measurement noise, and left the straggler
itself untouched. The next two subsections establish why, and the answer is
specific enough to act on.

\subsection{The benefit that does exist is coalescing, not skew handling}

AQE bundles several rewrites. To separate them we ran every AQE cell twice,
identical but for \texttt{spark.sql.adaptive.coalescePartitions.enabled}:

\begin{center}
\begin{tabular}{lrrr}
\toprule
$P$ & \PThin & \PMid & \PThick \\
\midrule
AQE, coalescing on (s)   & \AqeCoalWallThin & \AqeCoalWallMid & \AqeCoalWallThick \\
AQE, coalescing off (s)  & \AqeNoCoalWallThin & \AqeNoCoalWallMid & \AqeNoCoalWallThick \\
Baseline (s)             & \BaseCoalWallPreciseThin & \BaseCoalWallPreciseMid & \BaseCoalWallPreciseThick \\
\midrule
Longest task, coal.\ off (s) & \AqeNoCoalMaxTaskThin & \AqeNoCoalMaxTaskMid & \AqeNoCoalMaxTaskThick \\
\bottomrule
\end{tabular}
\end{center}

With coalescing disabled, AQE is indistinguishable from no AQE at all. The
longest task is unchanged in every column. Whatever small benefit AQE provided
came from merging small post-shuffle partitions --- which lowers fixed
per-task overhead and does precisely nothing for a straggler.

\subsection{The skew-join rule never fired, and why}

Reading the \emph{final} adaptive plan for every AQE cell explains the result
directly. Capturing it requires care: executing into a \texttt{noop} sink builds
a separate query execution and leaves the inspected plan reporting
\texttt{isFinalPlan=false}, showing the pre-adaptive shape and hiding exactly
the rewrite one wants to audit. Read after an action on the DataFrame itself,
every plan contains \texttt{AQEShuffleRead coalesced} and never
\texttt{AQEShuffleRead skewed}. \textbf{Spark's skew-join rule did not fire.}

The reason is a precondition that is documented but almost never discussed. A
skewed partition is split into pieces of roughly
\texttt{advisoryPartitionSizeInBytes}; if the hot partition is \emph{smaller}
than the advisory size, there is no split to make and the rule is a no-op
however skewed the partition is. The measured hot shuffle partition here is
\MeasHotMB~MB against a \NPART-partition advisory size of 16~MB.

This is a claim about a non-event, so it needs a positive control. We swept the
two documented trigger conditions to their most permissive settings ---
\texttt{skewedPartitionThresholdInBytes} from 64~MB down to 256~KB, and
\texttt{skewedPartitionFactor} from 5 down to 1. The rule still never fired.
Lowering only the advisory size engaged it immediately:

\begin{center}
\begin{tabular}{lrrl}
\toprule
Partitions & Advisory & Threshold & Skew-join applied \\
\midrule
16  & 16\,MB   & 8\,MB    & no  \\
16  & 1\,MB    & 1\,MB    & \textbf{yes} \\
64  & 1\,MB    & 1\,MB    & \textbf{yes} \\
200 & 256\,KB  & 256\,KB  & \textbf{yes} \\
\bottomrule
\end{tabular}
\end{center}

One configuration change, no code change. We regard this as the most
practically consequential finding in the paper: practitioners tune the two
parameters whose names contain the word \emph{skew}, and the parameter that
actually gates the optimisation at moderate data volumes is the one that does
not. Section~\ref{sec:decision}'s procedure encodes all three preconditions for
this reason, and predicted the non-event from data statistics alone.

\subsection{Salting worked, and its advantage grows with the probe side}

\begin{center}
\begin{tabular}{lrrr}
\toprule
$P$ & \PThin & \PMid & \PThick \\
\midrule
Best selective-salt $k$        & \SaltBestKThin & \SaltBestKMid & \SaltBestKThick \\
Salting speedup                & \SaltSpeedupThin$\times$ & \SaltSpeedupMid$\times$ & \SaltSpeedupThick$\times$ \\
Longest task, baseline (s)     & \BaseMaxTaskThin & \BaseMaxTaskMid & \BaseMaxTaskThick \\
Longest task, salted (s)       & \SaltMaxTaskThin & \SaltMaxTaskMid & \SaltMaxTaskThick \\
\textbf{Critical-path cut}     & \textbf{\SaltTaskCutThin$\times$} & \textbf{\SaltTaskCutMid$\times$} & \textbf{\SaltTaskCutThick$\times$} \\
\bottomrule
\end{tabular}
\end{center}

At $P=\PThin$ the result is inside the 18.3\% floor and we draw no conclusion.
At $P=\PMid$ and $P=\PThick$, where the floor is 0.1\% and 6.5\%, salting
returns \SaltSpeedupMid$\times$ and \SaltSpeedupThick$\times$ --- comfortably
outside noise --- and cuts the critical path by \SaltTaskCutMid$\times$ and
\SaltTaskCutThick$\times$ where AQE cut it by nothing.

One nuance worth recording, because it bears on the cost model. The $k$ that
minimises \emph{wall-clock} is \SaltBestKThick{} at $P=\PThick$, but a larger
$k$ minimises the \emph{critical path}: at $k=16$ the longest task falls
further still. The two optima differ because on two task slots the extra
shuffle from a larger $k$ costs more than the additional straggler relief
returns. On a cluster with more slots they should converge, and that is one of
the things the released cluster configuration is designed to check.

\subsection{Broadcast dominates at this scale, and we say so}

The broadcast arm completes in \BcastWallThin~s, \BcastWallMid~s and
\BcastWallThick~s --- faster than every shuffle-join strategy at every probe
thickness. This is expected and it is worth stating plainly rather than
omitting: our dimension table is well under Spark's default broadcast
threshold, and a practitioner in that regime should broadcast and stop reading.
Automatic broadcast is disabled throughout precisely so that the shuffle-join
behaviour we set out to study is not silently replaced. The results here apply
to joins whose dimension side exceeds the broadcast threshold, which is the
regime in which skew mitigation is a live question at all.

\subsection{Replication cost matches the model exactly}

Shuffle volume is a property of the plan and the data rather than of the
machine, so it is the measurement that transfers off a single node. Measured in
shuffle \emph{records} it does more than transfer: it matches
Equation~\eqref{eq:model} to the row.

An earlier version of this analysis measured replication in shuffle
\emph{megabytes} and reported a sevenfold gap between the two salting variants.
That was wrong, and the error is worth recording because it is easy to repeat.
Selective salting replicates only $P(k-1)$ rows --- a few hundred across the
entire sweep --- while the megabyte growth was dominated by the one-column
\texttt{\_salt} addition to every one of two million fact rows. Over 99\% of
what was attributed to replication was the extra column. Records have no such
confound: they are exactly the quantity the model's $bP(k-1)$ term refers to.

Fitting a least-squares line to shuffle records against $k$ over a matched
range gives:

\begin{center}
\begin{tabular}{lrrr}
\toprule
Variant & records per unit $k$ & model predicts & $R^{2}$ \\
\midrule
Selective salting & \SelRecSlope & $P = \PHOT$ & \SelRecR \\
Uniform salting   & \UnifRecSlope & $M = \NDIM$ & \UnifRecR \\
\bottomrule
\end{tabular}
\end{center}

Selective salting adds \SelRecSlope{} rows per salt value, and the hot key has
\PHOT{} rows on the probe side. Uniform salting adds \UnifRecSlope{} rows per
salt value, and the probe table has \NDIM{} rows. The measured ratio,
\RecSlopeRatio, equals $M/P$ exactly, which is what the model predicts, and
both fits are linear to $R^{2} = 1.000$.

This is the paper's cleanest quantitative result, and we would emphasise
\emph{why} it is clean rather than merely report it. Records are counted, not
timed. They do not vary between repetitions, they do not depend on the machine,
and they are not confounded by anything the scheduler does. The replication
term of the cost model is therefore not an approximation fitted to noisy
observations --- it is exact, and the distinction between the two salting
variants is a factor of $M/P$, which for any realistic fact--dimension join is
three orders of magnitude, not seven-fold.

The practical reading is that \textbf{the tutorial formulation of salting is
the expensive one}, and by a margin far wider than the folklore suggests.

\begin{table}[t]
\centering
\caption{Best configuration per strategy, with speedup over the unmitigated baseline. $\mu$ is the probe-side hot-key multiplier; skew is the join-stage maximum-to-median task duration ratio.}
\label{tab:headline}
\footnotesize
\begin{tabular}{lrrrrr}
\toprule
$\mu$ & strategy & $k$ & wall (s) & speedup & skew \\
\midrule
1.00 & baseline & 1 & 2.61 & 1.00 & 7.48 \\
1.00 & broadcast & 1 & 2.62 & 1.00 & 1.05 \\
1.00 & salt\_selective & 2 & 2.69 & 0.97 & 5.82 \\
1.00 & salt\_uniform & 2 & 2.78 & 0.94 & 4.97 \\
1.00 & aqe & 1 & 2.88 & 0.91 & 1.02 \\
1.00 & aqe\_salt & 8 & 3.01 & 0.87 & 2.50 \\
10.00 & salt\_selective & 2 & 5.19 & 1.34 & 23.87 \\
10.00 & aqe\_salt & 8 & 5.53 & 1.26 & 1.32 \\
10.00 & broadcast & 1 & 5.60 & 1.24 & 1.03 \\
10.00 & salt\_uniform & 8 & 5.68 & 1.23 & 6.39 \\
10.00 & aqe & 1 & 6.90 & 1.01 & 3.54 \\
10.00 & baseline & 1 & 6.96 & 1.00 & 35.97 \\
32.00 & salt\_selective & 2 & 11.08 & 1.68 & 68.38 \\
32.00 & aqe\_salt & 8 & 11.67 & 1.60 & 13.91 \\
32.00 & salt\_uniform & 8 & 11.72 & 1.59 & 17.45 \\
32.00 & broadcast & 1 & 13.32 & 1.40 & 1.03 \\
32.00 & aqe & 1 & 18.33 & 1.02 & 10.39 \\
32.00 & baseline & 1 & 18.62 & 1.00 & 101.01 \\
\bottomrule
\end{tabular}
\end{table}


\subsection{Where the cost model held, and where it did not}

The model's structural claims are supported. The replication term is exact, as
above. Salting monotonically reduces the critical path: at the thickest probe
side the longest join-stage task falls from \BaseWallThick~s unmitigated to a
small fraction of that under salting.

The \emph{convexity} of $T(k)$ in wall-clock, however, did not resolve, and we
report that plainly rather than presenting a fitted curve as a confirmation.

Fitting Equation~\eqref{eq:model} to the selective-salting sweeps gives
$R^{2}$ of \RsqThin, \RsqMid{} and \RsqThick. Those look respectable and are
misleading. The quantity that matters is whether the replication coefficient
$b$ --- the entire denominator of $k^{*}$ --- is resolved by the data at all:

\begin{center}
\begin{tabular}{lrrr}
\toprule
$P$ & \PThin & \PMid & \PThick \\
\midrule
$R^{2}$ & \RsqThin & \RsqMid & \RsqThick \\
$|b| / \mathrm{se}(b)$ & \BtstatThin & \BtstatMid & \BtstatThick \\
$b$ identified? & \IdentThin & \IdentMid & \IdentThick \\
$k^{*}$ & \KstarThin & \KstarMid & \KstarThick \\
\bottomrule
\end{tabular}
\end{center}

At two of three workloads $b$ is statistically indistinguishable from zero. A
$k^{*}$ computed from such a fit is an artefact of noise, and our
implementation refuses to return one: \texttt{optimal\_k\_unconstrained} yields
\textsc{NaN} rather than a number, and the recommendation falls back to the
saturation and parallelism ceilings alone. Only \NumIdentified{} of
\NumFits{} fits supports a cost-optimum at all.

A second, independent sign of trouble: the fitted $\gamma = a/b$ varies by
\GammaSpread$\times$ across three workloads that differ only in $P$, on one
machine, in one session. Section~\ref{sec:costmodel} advances $\gamma$-stability
as the property that would make the closed form practically useful. Our own
data contradict it.

We therefore make no claim that the closed form is validated. What we can say
is precise: the replication \emph{term} is exact; the replication \emph{cost in
time} is, at this scale, too small for the convexity it implies to be
observable; and the conditions under which it should become observable ---
cluster-scale network shuffle, or a substantially thicker probe side --- are
identified and configured in the released artifact.

The practical consequence points the opposite way from the folklore, and is
worth stating for practitioners rather than only for reviewers. \textbf{With
selective salting, over-provisioning $k$ is close to free.} The warning against
large $k$ is well founded for the uniform variant, where the replication cost
is $M/P$ times larger and clearly measurable. For the selective variant, the
sensible advice is to salt at the saturation bound and stop worrying about the
exact optimum.

\subsection{Summary}

\begin{itemize}
\item AQE's advantage over an unmitigated baseline fell from
\AqeSpeedupThin$\times$ to \AqeSpeedupThick$\times$ as the probe-side hot key
grew from \PThin{} to \PThick{} rows, while salting's rose from
\SaltSpeedupThin$\times$ to \SaltSpeedupThick$\times$.
\item Spark's skew-join rule never fired on a join with a
\BaseSkewThick:1 task-time skew ratio, verified from final adaptive plans. The
measured AQE benefit was partition coalescing.
\item Neither documented trigger knob was responsible; the binding condition
was \texttt{advisoryPartitionSizeInBytes} exceeding the hot partition. Lowering
it engaged the rule immediately.
\item Selective salting replicates \SelRecSlope{} rows per salt value against
uniform salting's \UnifRecSlope{} --- a ratio of \RecSlopeRatio, exactly the
$M/P$ the cost model predicts, with both fits linear to $R^{2}=1.000$.
\item The replication term of the cost model is confirmed exactly. Its
consequence for runtime is not: the coefficient $b$ is statistically
unidentified at two of three workloads, fitted $\gamma$ varies by
\GammaSpread$\times$ on one machine, and we therefore do not claim the closed
form for $k^{*}$ is validated.
\end{itemize}

\section{Threats to Validity}
\label{sec:threats}

\subsection{Single-node measurement}

This is the principal threat and we state it without hedging. All measurements
were taken on one machine with a small number of cores. Absolute wall-clock
times therefore do not transfer to a cluster, and we make no claim that they
do.

What survives the objection is specific. Shuffle volume and spill volume are
determined by the physical plan and the data partitioning, not by the machine:
selective salting with $k=32$ replicates the same probe-side rows on a laptop
as on a hundred nodes. The skew ratio is a within-run ratio of task durations
and cancels most machine-specific scaling. And the model's structure ---
convexity in $k$, the $\sqrt{H/P}$ dependence, the two ceilings --- is
analytic, not measured; the experiments test whether it holds, and the
calibration constant $\gamma$ is expected to differ per platform, which is why
it is a calibrated parameter and not a published number.

What does \emph{not} survive is any claim about network shuffle. A single-node
Spark shuffle writes to local disk and reads back; a cluster shuffle crosses the
network, where the replication cost of salting is materially higher. Our
estimate of $b$ is therefore \emph{conservative}, and the true $k^{*}$ on a
cluster should be \emph{smaller} than we report --- which strengthens rather
than weakens the paper's central practical claim that customary salt
cardinalities are too large. We flag this asymmetry explicitly so that a reader
does not have to infer it.

The parallelism ceiling is the sharpest limitation. With few task slots, $C$
binds in Equation~\eqref{eq:keff} for most configurations, so our measurements
exercise that branch more than the cost-optimum branch. We address this by
reporting the model's projection across $C$ rather than only the measured
point, and by being clear that the projection is a projection.

\subsection{Synthetic data}

Zipf is a model of skew, not skew itself. Real key distributions have
structure --- correlated hot keys, temporal drift, keys hot on one side only in
certain periods --- that a stationary Zipf law does not reproduce. We mitigate
this by using an established distributional form, by sweeping $\theta$ rather
than fixing it, by adopting a schema from a real prognostics
dataset~\cite{saxena2008cmapss}, and by including a multi-hot-key
configuration. The harness accepts skewed TPC-H~\cite{microsoft_skewgen} and
real C-MAPSS data as drop-in replacements. But a study on production key
distributions would be stronger, and we do not have one.

\subsection{Bounded probe-side multiplicity}

Our two-sided workloads cap the hot key's probe-side representation at $\mu$
times the mean rather than letting it follow the fact side's Zipf law. This is
a deliberate bound, justified in Section~\ref{sec:methodology}, but it is a
restriction on generality: production joins do exist in which both sides are
severely concentrated and the output genuinely explodes. For those, the binding
problem is output cardinality rather than straggler relief, and neither
mechanism studied here is the right answer --- the correct response is to
change the query, not to tune the join. We regard that as a different problem,
and we do not claim our results speak to it.

\subsection{Version specificity}

AQE's skew-join rule, its defaults, and the interaction between
\texttt{coalescePartitions} and \texttt{skewJoin} all change between Spark
releases, and the production descendant described by Xue et
al.~\cite{xue2024adaptive} differs from the open-source rule. Our measurements
pin one version, recorded in the artifact. The cost model does not depend on
the version; the decision procedure's Branch 2, which reasons about the trigger
condition, does.

\subsection{Construct validity of the skew ratio}

Maximum over median task duration is a proxy, and it has a confound serious
enough that we no longer use it as a primary metric.

The problem is the denominator. AQE's partition coalescing merges the sixteen
shuffle partitions into three, which changes what ``the median task'' refers to
without changing the straggler at all. Comparing a three-task stage's
max-over-median against a sixteen-task stage's is comparing different
quantities. In our data this inverts the ranking: at the thickest probe side
the AQE arm reports a far lower skew ratio than the unmitigated baseline while
its slowest task is \emph{longer} in absolute terms. Any reader who took the
ratio at face value would conclude the opposite of the truth.

We therefore report three things instead. The \emph{absolute maximum task
duration} in the join stage, which is the critical path and has no denominator
to move. The \emph{straggler index}, maximum over \emph{mean} task duration:
because total task time is roughly conserved under coalescing, the mean is
stable where the median is not. And \emph{task counts} alongside every ratio,
so that a reader can see immediately when two numbers are not comparable. The
max-over-median ratio is retained only for continuity with the practitioner
literature, which universally quotes it.

The residual confounds remain: a near-zero median makes the ratio unstable, and
task-duration variance from scheduling and garbage collection is not skew. Both
are reasons to lead with shuffle-byte and record distributions, which have no
such confound.

\subsection{The dimension side is small enough to broadcast}

Our probe table is well under Spark's default ten-megabyte automatic broadcast
threshold. In a default configuration this join would be planned as a broadcast
hash join, which has no shuffle and therefore no shuffle skew, and neither
mechanism we study would be engaged at all. We disable automatic broadcast
throughout precisely to isolate the shuffle-join behaviour we set out to
measure, and we run the broadcast arm so that its dominance at this scale is
visible in the results rather than hidden by the configuration.

The honest scope statement follows: our results apply to fact--dimension joins
in which the dimension side exceeds the broadcast threshold, which is the
regime where skew mitigation is a live question. A reader whose dimension table
fits in a broadcast should broadcast it and stop reading.

\subsection{Execution order and the noise floor}

An earlier version of this harness executed cells in configuration order, which
placed the unmitigated baseline first in every workload block --- on a cold
JVM, and therefore systematically slow. Because every speedup is measured
against that baseline, the bias inflated all of them.

We detected this with a control the harness generates for free: selective
salting at $k=1$ compiles to exactly the same physical plan as the baseline, so
any difference between the two is measurement noise. The discrepancy reached
17.7\% on the first workload while remaining under 0.3\% on the later ones ---
a clear order effect rather than random variation. Cell order is now
randomised under a fixed seed, a session-level warmup precedes the grid, and
the control is computed and reported for every run. It is included in the
artifact.

We state the consequence rather than burying it: \textbf{no speedup smaller
than the control discrepancy should be read as a result}, in this paper or in
any run of this harness.

\subsection{Statistical power}

Five repetitions per cell, and three for the configuration probes, is few. We
report medians and interquartile ranges rather than means and standard
deviations, and we do not attach $p$-values to five-against-five comparisons,
where the smallest attainable two-sided rank-test $p$ is above the conventional
threshold in any case. Differences of a few percent in wall-clock are within
noise at this sample size and we do not interpret them. The deterministic
counters --- shuffle bytes and records --- carry the claims that need
precision, and they do not vary between repetitions at all.

\subsection{One partition count, one core count}

Every reported cell uses sixteen shuffle partitions and two task slots. The
parallelism ceiling therefore binds the recommendation in essentially every
configuration, which means the cost-optimum branch of
Equation~\eqref{eq:keff} --- the branch the closed form exists to serve --- is
never the operative one in our own data. This is the sharpest reason the
cluster validation described in Section~\ref{sec:conclusion} matters, and it
is why we release the cluster configuration alongside the single-node one.

\subsection{Generalisation beyond inner equi-joins}

We evaluate inner equi-joins on a single key. Outer joins, multi-key joins,
joins whose skewed key is also the aggregation key, and joins feeding a
subsequent shuffle all plausibly behave differently. The model's derivation
does not obviously break for these cases, but we have not tested them, and we
do not claim them.

\subsection{Absence claims}

Section~\ref{sec:background} rests on two negative results: no independent
evaluation of Spark's skew-join rule, and no peer-reviewed treatment of
salting. Negative results in literature search are only as good as the search.
Ours covered dblp, the arXiv full-metadata API, Crossref, and domain-scoped
searches across the major publishers, but Semantic Scholar and OpenAlex were
unavailable, and at least one relevant paper~\cite{phan2022skewjoin} is not
indexed in dblp. We therefore present these as strong evidence of a gap rather
than as proof of one, and we cite the nearest existing work explicitly rather
than claiming an empty field.

\section{Conclusion}
\label{sec:conclusion}

Spark ships an automatic remedy for skewed joins, and practitioners continue to
salt by hand. We set out to determine which they should use, and how much salt
to use when salting. The answer on our workload was blunter than expected:
Spark's skew-join rule never engaged at all, so there was nothing to choose
between.

The measurement is a null result and we report it as one. Enabling AQE on a
join whose longest task ran \BaseMaxTaskThick~s produced a longest task of
\AqeCoalMaxTaskThick~s, with speedups of \AqeSpeedupMid$\times$ and
\AqeSpeedupThick$\times$ against measurement floors of 0.1\% and 6.5\%. A
paired control with coalescing disabled shows that what little AQE contributed
was partition merging, which lowers fixed overhead and does nothing for a
straggler. Selective salting, on the same workload, cut the critical path
\SaltTaskCutThick$\times$.

The cause is a precondition that is documented and almost never discussed. A
skewed partition is split into pieces of \texttt{advisoryPartitionSizeInBytes};
a hot partition smaller than that size cannot be split at all, and the rule
becomes a no-op however extreme the skew. Our hot partition measured
\MeasHotMB~MB against a 16~MB advisory size. Sweeping the two parameters whose
names contain the word \emph{skew} to their most permissive settings changed
nothing; lowering the advisory size engaged the rule immediately. Practitioners
tune the two knobs that are named for the problem, and the one that gates the
solution at moderate data volumes is not among them.

On the cost model we are more circumspect than we expected to be. The
replication term is confirmed exactly in shuffle records --- \SelRecSlope{}
rows per salt value for selective salting against \UnifRecSlope{} for uniform,
a ratio of \RecSlopeRatio{} matching the predicted $M/P$ with $R^{2}=1.000$.
Its consequence for runtime is not confirmed: the replication coefficient is
statistically unidentified in all \NumFits{} fits, and in \NumNegativeB{} of
them it comes out with the wrong sign. At this scale replication cost is not
merely unresolved in wall-clock; it is absent. We therefore present the closed
form $k^{*}=\sqrt{\gamma H/P}$ as theory, state plainly that this study does
not validate it, and release the cluster configuration under which the
prediction --- that network shuffle raises the coefficient and lowers $k^{*}$
--- becomes testable.

The decision between mechanisms turns on a statistic that costs one sampled
group-by to obtain. When skew is one-sided and above AQE's trigger, AQE alone
suffices and manual salting adds replication for little gain. When the same key
is heavy on both sides, split-and-replicate degenerates and salting remains
necessary. When the hot partition sits below AQE's absolute threshold, the rule
does nothing at all --- silently --- and a practitioner who has enabled AQE
believes the problem is handled.

The artifact matters as much as the result. We cannot out-scale the
organisations whose systems papers describe this machinery at production
volume, and we do not try. What a practitioner-built harness can contribute is
specification: a public, reproducible generator with a declared skew taxonomy,
a metrics layer that separates deterministic counters from machine-dependent
timings, and an honest account of what a single node does and does not license.
\textsc{skewbench} is released under the MIT licence.

Three directions follow. The most useful is validation at cluster scale, where
the network shuffle makes $b$ larger and $k^{*}$ correspondingly smaller than
we report; our prediction is directional and testable. The second is to fold
$\gamma$ estimation into the engine, so that $k^{*}$ is computed rather than
configured --- salting is currently the one skew remedy Spark cannot apply for
you, and the model here is simple enough to be a planner rule. The third is to
extend the taxonomy: multi-key joins, outer joins, and skew that drifts over
time are all common in production and none is covered here.


\section*{Artifact Availability}
The \textsc{skewbench} harness, the workload generator, the raw event logs, the
analysis pipeline, and the scripts that produce every figure and table in this
paper are released under the MIT licence and archived with a DOI. All results
reported here are reproducible with a single command on a commodity laptop.

\bibliographystyle{IEEEtran}
\bibliography{refs}

\end{document}
